% SPDX-License-Identifier: CC-BY-4.0
% =====================================================================
% Appendix A of D. Joyner and Claude, "A hitch-hiker's guide to the
% metaplectic trace formula for SL(2)" (November 2026).
%
% Comparisons with the 1995 notes and the 1997 book manuscript. Until the
% second referee report (ChatGPT, November 2026) these paragraphs stood in
% the body of Sections 2-6, after the statements they concern; they were
% moved here unchanged except for the lead-ins. The numbering maps
% (1995 -> 2026) are in ../revision_log.md. The subsection on Section 7 was
% added in stage 6.
% =====================================================================

\section{Comparison with the 1995 notes}\label{sec:appendix-1995}

Sections~\ref{sec:background}--\ref{sec:stable} revise Sections~2--7
of the 1995 notes \emph{A hitch-hiker's guide to the metaplectic trace
formula for $SL(2)$: basic theory} (D.~Joyner, January 1995, revised June
2002), and the parts of Chapter~9 of the 1997 book manuscript \emph{An
application of Arthur's invariant trace formula for $SL(2)$} from which
they came. This appendix records, statement by statement, what the
revision replaces and what it corrects. File names such as
\texttt{book7d.tex} refer to that manuscript. Item numbers such as
``II.5'' or ``A.2'' refer to the first referee report on the 1995 notes
(Claude, November 2026); the list of items, and the tables matching the
1995 numbers with the present ones, are in the file
\texttt{revision\_log.md} that accompanies the source. Nothing in this
appendix is used in the body of the paper.


\subsection{Section~\ref{sec:background}: the cover}

\begin{description}
\item[Section~\ref{sec:local}, the cover of $G(\Q_p)$.] The 1995 notes call $\overline{G(\Q_p)}$ a subgroup of $\prod_{v\mid p}\overline{SL(2,k_v)}$; it is a quotient.

\item[The height $H$, \eqref{eq:H}.] The 1995 notes define $H$ with the maximum of the $|t_v|_v$, $v\mid p$, in place of their product. The product is forced by $\langle H(x),\chi_1\rangle=\log|\chi_1(\tilde a)|$ (Section~\ref{sec:chambers}).

\item[Section~\ref{sec:chambers}.] The 1995 notes call the fundamental coweight $\varpi$ the dual root. An earlier draft of this revision had $\alpha=2\chi_1$ and $X_\alpha=1$, which is correct only for $k=\Q$ (Remark~\ref{rem:alpha}).

\item[Section~\ref{sec:functions}, genuine and anti-genuine.] The 1995 notes use genuine functions in~\S2.1.4 and $\pi(g,\zeta)=\zeta^{-1}\pi(g,1)$ in~\S3.1. For $n=2$ both conventions agree with Li's; for general $n$ they do not (Remark~\ref{rem:general-n}).

\item[Section~\ref{sec:functions}, the unit element.] The 1995 notes identify $K_p$ with $\{(k,1)\}$, which suggests the function $(k,\zeta)\mapsto\zeta^{-1}$ as the unit of the Hecke algebra; the unit is $e_p$ of~\eqref{eq:unit}. An earlier draft of this revision defined $C_c^\infty(\overline{G(\Q_p)})_-$ by requiring $g\mapsto f(g,1)$ to be smooth; that definition excludes $e_p$ and was replaced after the 2026 referee report by ChatGPT.

\item[The integration formula~\eqref{eq:iwasawa}.] The 1995 notes have $\delta_B(a)^{1/2}$ for $g=kan$; the exponent must be $1$ for $g=kan$ and $-1$ for $g=nak$. In Section~4 they have $\delta(a)$ for $x=nak$.

\item[Lemma~\ref{lem:global}(a), the splitting over $G(\Q)$.] The 1995 notes define the splitting as $\gamma\mapsto(\gamma,
s_\A(\gamma)^{-1})$, with $s_v(\gamma)^{-1}=(c,d)_v$ when $v$ is finite,
$cd\ne0$ and $n\nmid\ord_v(c)$, and with the cocycle $\prod_v\beta_v$.
That product does not converge, and the formula belongs to a different
normalization. In the normalization~\eqref{eq:sigma-A},
$\kappa_\A$ of~\eqref{eq:kappa-A} replaces $s_\A^{-1}$. For $n=2$ the two
agree at every non-dyadic $v$ with $\gamma\in K_v$: if $\ord_v(c)>0$ is
even, or if $\ord_v(c)=0$, both $(c,d)_v$ and $\kappa_v(\gamma)$ are $1$.
They differ only at the finitely many places where $\gamma\notin K_v$,
and at dyadic places.

\end{description}

\subsection{Section~\ref{sec:rank0}: rank 0}

\begin{description}
\item[The cocycle on $A$, \eqref{eq:cocycle-A}.] The 1995 formula $\beta(g_1,g_2)=(g_1^{-1},g_2^{-1})(-g_1g_2^{-1},g_1^{-1}g_2^{-1})$ on $A$ reduces to~\eqref{eq:cocycle-A} for $n=2$.

\item[Lemma~\ref{lem:genuine-chars}.] This replaces Lemma~3.1 of the 1995 notes, which cited~\cite{Ge2} for factorizability of genuine representations of the restricted product. For $n=2$ the statement concerns the quotient and needs no citation.

\item[Theorem~\ref{thm:corr}.] This is the $n=2$ form of Theorem~3.2 of the 1995 notes, restated as a correspondence of automorphic spectra. The 1995 statement matches genuine representations with characters $\chi^N$, $N=n/\gcd(n,2)$; for $n=2$, $N=1$, and the matching requires the choice of $\bar\chi_0$, which the 1995 text did not make.

\item[Theorem~\ref{thm:tf-rank0}.] This replaces Theorems~3.8 and~3.9 of the 1995 notes, whose orbital integrals reduce to the values $f(\delta)$ for $n=2$.

\item[Lemma~\ref{lem:transfer}.] This replaces Lemmas~3.11 and~3.12 and the argument through Claims~1 and~2 of the 1995 notes. The local maps $\bar f_v\mapsto f_v$ involve no sum over $T(\Q_p)_\delta\backslash T(\Q_p)$ and no factor $s_\A(\delta)$, because there is no conjugation.

\end{description}

\subsection{Section~\ref{sec:geometric}: the geometric side}

\begin{description}
\item[Lemma~\ref{lem:classes}.] The 1995 text lists the centralizer $G$ for $\gamma=\pm1$ among the centralizers of semisimple classes that are ``anisotropic modulo the center'' and calls all such classes unramified. Neither holds for $\pm1$: $G$ is isotropic, and $\mathcal O_{\pm1}$ contains $\pm$(unipotent) elements. The vanishing of genuine orbital integrals for general $n$ is \texttt{book3b.tex}, Lemma~5.2.

\item[Lemma~\ref{lem:kernels}.] This is Lemma~4.1 of the 1995 notes with $K_P=\sum_{\mathcal O}
K_{P,\mathcal O}$ (not $\sum K_{\mathcal O}$) and with one kernel
convention, $f(x^{-1}\gamma y)$, throughout; no passage to the
archimedean place is needed for the convergence.

\item[Lemma~\ref{lem:height-ineq}.] Part~(b) replaces Lemma~4.2 of the 1995 notes, which has an unused $N$ and only the bound $\le1$, cited from~\cite{A2}. The constants $u\ge1/2$ and $T_0>\sqrt3/2$ of the 1995 text are the classical ones for $SL(2,\Z)$ over $\Q$ and are not used.

\item[Lemma~\ref{lem:partition}.] This is Lemma~4.3 and (4.2) of the 1995 notes, there with two parameters $T_0>T_1$ in reversed roles and with a function $F^B$ that is identically $1$ for $SL(2)$.

\item[Theorem~\ref{thm:coarse}.] The theorem replaces Lemma~4.4, Corollaries~4.5--4.7, Proposition~4.8 and Corollary~4.9 of the 1995 notes. Their proof of Lemma~4.4 omits the terms $\gamma\notin B(\Q)$ and the sum over $\mu\in A(\Q)$, and the per-class bounds in Corollaries~4.5--4.7 are not established there. Proposition~4.8 claims polynomiality for $K^T_{\mathcal O}$ and $J^T_{\mathcal O}$ separately, without a proof. Part~(b) of the theorem shows that for this group both integrals exist and coincide, and part~(d) gives the polynomials.

\item[Proposition~\ref{prop:hyperbolic}.] The 1995 weight $v_{\mathcal O}(x,T)=2(T-H(x))$ (item II.5 of the first referee report) is not left $A(\A)$-invariant; the $w$-term was lost in passing from the sum over $\delta$ in (4.3) to $\hat\tau_B(H(\eta x)-T)$ in (4.4). The weight $2T-\log\delta(x)-\log\delta(wx)$ of \texttt{book7d.tex} has the right shape but the wrong normalization, since $\log\delta_B=2H$ (item A.4).

\item[Proposition~\ref{prop:central}.] The 1995 formula (4.5) keeps only the $f(\mu)$-term of a central class (item II.6). In \texttt{book7d.tex}, $J_{\mathrm{unip}}$ is defined by truncating at $T_0=0$, which is correct here only because the integral converges at $T=0$; in general it must be the polynomial evaluated at $0$ (item A.5).

\item[Section~\ref{sec:local-dist}, test functions.] \texttt{book7d.tex} takes $f_\infty$ in $C^\infty$ rather than $C_c^\infty$, and asserts that Arthur's spaces of almost compactly supported and moderate functions coincide with $\mathcal H$ for both $G$ and $A$; for $A$ this fails.

\item[Lemma~\ref{lem:weight-local}.] At a $p$-adic place the weight is $\log\max(1,|y|_p)$, not
$\log\sqrt{|1+y^2|_p}$ as in \texttt{book7d.tex} and (7.13) of the 1995
notes (items II.22, A.6); the latter is the real formula.

\item[Lemma~\ref{lem:JA-central}.] The Example in \texttt{book7d.tex} sets $r_A=1$ and takes the plain limit of $J_A(az,f)$, which diverges (item A.2).

\item[Theorem~\ref{thm:fine}.] This is the corrected form of the Proposition of \texttt{book7d.tex} (cited in the 1995 Section~6 as ``Proposition~4.18''), which credits the result for covers to~\cite[Thm.~8.1]{A12}; Arthur proves it for connected reductive groups, and for covers it is Li's theorem~\cite{LiI} (item A.1). The sums there run over the rational class $\mathcal O$ instead of the $(G,S)$-classes (item A.7). In the 1995 text, 5.4(a) has only $a^G(S,z)f(z)+\tfrac12a^A(S,z)J_A(z,f)$; Theorems~\ref{thm:tf}(a) and~\ref{thm:itf}(a) use Theorem~\ref{thm:fine}(a) in full.

\end{description}

\subsection{Section~\ref{sec:spectral}: the spectral side}

\begin{description}
\item[Lemma~\ref{lem:spec-decomp}.] Lemma~5.1 of the 1995 notes asserts that the right regular representation
on $L^2(\Gamma\backslash G)$ decomposes into irreducibles ``each occurring
with finite multiplicity'' for every discrete $\Gamma$ in a unimodular
$G$. That is false in general (take $\Gamma=1$); what holds is a direct
integral decomposition for $G$ of type~I, and finite multiplicities in
the discrete part for lattices in reductive groups, as above.

\item[Lemma~\ref{lem:factor}.] This corrects Example~5.2 of the 1995 notes, which says that all genuine
unitary representations are factorizable (it should say irreducible),
that $\pi_v$ is unramified when it has a $\overline{K_v}$-fixed vector
(genuine representations have none; the condition is a $K_v^*$-fixed
vector), and which treats $\GA$ as a restricted product rather than a
quotient of one (item~I.4).

\item[Cuspidal data.] This replaces Example~5.4 of the 1995 notes, which says that two
parabolics with Levi $A$ are associate only if they are opposite (but $B$
is associate to itself), and that $(A,\rho)\sim(A,\rho')$ only if
$\rho=\rho'^w$; the correct condition is $\rho=\rho'$ or $\rho=\rho'^w$
(item~II.13).

\item[Proposition~\ref{prop:Lchi}.] The proposition replaces~(5.2)--(5.4), Corollary~5.5 and Proposition~5.6 of the 1995 notes. The 1995 description of $L^2(M(\Q)\backslash\overline{M(\A)^1})_\chi$ through the spaces $\Psi(P_1,P,\chi)$ is Langlands' original one; for $SL(2)$ the description in Proposition~\ref{prop:Lchi} is all that is used.

\item[The space $\mathcal H_B(\mu)$, \eqref{eq:HBmu}.] In Example~5.18 of the 1995 notes the condition $\varphi(ag)=\pi(a)\varphi(g)$ is imposed for all $a\in\overline{A(\A)}$, so that the parameter $\lambda$ is part of $\pi$ and the factor $\delta_B^{1/2}$ is missing.

\item[Theorem~\ref{thm:eisenstein}.] This is Theorem~5.8 of the 1995 notes, specialized to $P=B$, with three corrections. The measure on $i\mathfrak a^*$ is $dt/2\pi$ (the 1995 text has $(1/2\pi i)^{\dim A}$ in~(b) and $(1/2\pi)^{\dim A}$ in the norm of $L^\wedge_{\mathcal P}$). The number $n(A_B)$ of chambers in $\mathfrak a$ is $2$, not $1$. And the theorem is stated for one $\mu$, the spectral decomposition being Proposition~\ref{prop:Lchi}.

\item[Lemma~\ref{lem:intertwining}.] The 1995 notes leave $M(s,\pi)$ unnormalized (item~II.16).

\item[Theorem~\ref{thm:residual}.] The theorem corrects the Introduction of the 1995 notes, which says that the formula has no term corresponding to the residual term~\cite[(5.5)]{LL} (item~II.9).

\item[Proposition~\ref{prop:kernels-spec}.] This replaces Lemma~5.9, Theorem~5.11 and Examples~5.17--5.22 of the 1995
notes. Examples~5.21--5.22 there put the term
$I_B(\pi,f)\phi(x)\delta_B(x)\overline{\phi(y)}$ into $K_{G,\chi}$; it
belongs to $K_{B,\chi}$, and for $\chi$ cuspidal there is no Eisenstein
series in $K_{G,\chi}$ (item~II.13). Example~5.20 omits $\lambda$ and the
measure. The 1995 definition of $k^T_\chi$ has $K_{P,\chi}(x,x)$ where
$K_{P,\chi}(\delta x,\delta x)$ is meant (as in \texttt{book7f.tex}).

\item[Lemma~\ref{lem:truncation}.] This replaces Lemma~5.13 of the 1995 notes (proved there by reference to
the non-metaplectic case) and makes the range of $T$ explicit.
Lemma~5.14 and Example~5.12 there are the case $P=G$ of~\cite[Lemma~1.5]
{A3}; for $SL(2)$ that identity is~\eqref{eq:trunc} itself, since
$\hat\tau_B=\tau_B$, and we do not repeat it.

\item[Lemma~\ref{lem:kchi}.] This is Lemma~5.15 and Example~5.16 of the 1995 notes; their reference
``lemma~4.7'' for the partition is Lemma~\ref{lem:partition}.

\item[Theorem~\ref{thm:Jchi-poly}.] The theorem replaces Lemmas~5.23, 5.25 and~5.26 of the 1995 notes, whose proofs are incomplete (item~II.12). The proof of Lemma~5.23 asserts that $\Lambda^T_2K_\chi(x,x)$ is compactly supported; it is only rapidly decreasing (Lemma~\ref{lem:truncation}(c)). The proof of Lemma~5.25 (``$K_P$ is continuous, hence integrable'') does not give integrability. And the equality $K^T_\chi=J^T_\chi$ claimed there, of $\int\Lambda^T_2K_\chi$ and $\int\sum E\,\overline{\Lambda^TE}$, is a consequence of Lemma~\ref{lem:truncation}(b), not the content of Arthur's~\cite[Lemma~2.4]{A3}; what is needed is~(c), which compares them with $\int k^T_\chi$.

\item[Theorem~\ref{thm:Jchi}.] Theorem~\ref{thm:Jchi} replaces Theorem~5.27 and Example~5.28 of the
1995 notes. Three things change.
\begin{itemize}
\item \emph{The factor of $T$} (item~II.11). Arthur's formula
  \[
  M_B^T=\mathrm{vol}(\mathfrak a/\Z\alpha^\vee)\lim_{\zeta\to0}
  \sum_{s\in\{1,w\}}\frac{e^{(s\zeta)(T)}M(s,\lambda)^{-1}
  M(s,\lambda+\zeta)}{\theta(s\zeta)},
  \]
  with $\theta(\zeta)=\zeta(\alpha^\vee)=d\zeta$
  and $\mathrm{vol}(\mathfrak a/\Z\alpha^\vee)=d$ for our $dX$
  (Remark~\ref{rem:alpha}), gives~\eqref{eq:MT} for every $k$; the factor
  $d$ cancels, so the correction $\alpha^\vee=d$ (Remark~\ref{rem:alpha}) does not change
  $M_B^T$ once the volume is taken with the same measure. The 1995
  computation identifies $\varsigma(T)=2\varsigma T$ and
  $\varsigma(\alpha^\vee)=2\varsigma$ (correct for $k=\Q$, where
  $\lambda=2\varsigma$), takes the volume to be $1$, and then omits the
  division by $\varsigma(\alpha^\vee)$; the limit is
  $2T-\frac12M^{-1}\frac{dM}{d\varsigma}$, which is~\eqref{eq:MT}, not
  $4T-M^{-1}M'$. The sign is right, as the 1995 text suspected and as
  in~\cite{LL}; Proposition~\ref{prop:Tcoef} confirms the factor.
\item \emph{Ramified classes} (item~II.10). The 1995 Theorem~5.27 assumes
  $\chi$ unramified, and its Example~5.28 says correctly that $(A,\mu)$
  is unramified if and only if $\mu^w\ne\mu$, but the classes with
  $\mu^w=\mu$ never appear. By Lemma~\ref{lem:w-action}(d) they are the
  $\chi_b$, and they carry the residual spectrum and the term
  $\frac14\operatorname{tr}(M(w,0)I(\mu_0,f))$. In Arthur's fine
  expansion~\cite[Thm.~8.2]{A14}, \cite{LiIII} this is the term
  $|W^A_0||W^G_0|^{-1}|\det(w-1)_{\mathfrak a}|^{-1}
  \operatorname{tr}(M(w,0)I(\mu_0,f))$, with $|W^A_0||W^G_0|^{-1}=\frac12$
  and $|\det(w-1)_{\mathfrak a}|=2$. In the trace formula for $GL(2)$
  of~\cite[\S6]{GJ} the same term is written $-\frac14\operatorname{tr}
  (M(0)\pi_0(f))$, next to $+\frac1{4\pi}\int\operatorname{tr}(M(-iy)M'(iy)
  \pi_{iy}(f))\,dy$: that formula has the trace of $R(f)$ on the discrete
  spectrum alone on one side, and moving both terms to the spectral side
  gives the signs of Theorem~\ref{thm:tf}(e). The sign of the cross terms
  in~\eqref{eq:MS}, from which it comes, is also checked numerically on
  the Eisenstein series of $SL(2,\Z)$ (Remark~\ref{rem:verified-sec6}),
  and $M(w,0)$ is computed in Proposition~\ref{prop:Mw0}.
\item The lines ``$M_G^T(\pi)_\chi=1$'' and ``($P=B$, $\chi=(G,\rho)$)
  $M_B^T(\pi)_\chi=1$'' of Example~5.28 are dropped: for cuspidal $\chi$
  there is no $B$-term.
\end{itemize}

\item[Theorem~\ref{thm:tf}.] In the 1995 notes, Section~5.4 lists~(a) as
$a^G(S,z)f(z)+\frac12a^A(S,z)J_A(z,f)$, without the unipotent terms
$a^G(S,u_\xi)J_G(zu_\xi,f_S)$; \texttt{book7f.tex} has them in the form
$\sum_{\gamma\in\mathcal O}a^G(S,\gamma)J_G(\gamma,f)$, a sum over the
rational class instead of the $(G,S)$-classes (items~II.6, A.7).
Item~(e) there is $\frac12\int\operatorname{tr}(M_A(\pi)_\chi I_B(\pi,f)_
\chi)\,d\pi$ over $\chi\in\mathcal X\smallsetminus\mathcal X_G$ with
$M_A=M_A^T|_{T=0}$, without the residual and ramified terms. Section~5.4
also uses $S$ and $\A_S$ before Section~6 defines them; they are now
defined in Section~\ref{sec:local-dist}, and ``given in section~5''
refers to Section~\ref{sec:geometric}.

\end{description}

\subsection{Section~\ref{sec:invariant}: the invariant form}

\begin{description}
\item[The introduction of Section~\ref{sec:invariant}.] Compared with Section~6 of the 1995 notes:
\begin{itemize}
\item Conjectures~6.3, 6.4, 6.7, 6.8 and~6.10 there (the trace
  Paley--Wiener theorems and the normalization of intertwining
  operators) are no longer hypotheses: 6.10 is a theorem of
  Li~\cite{LiII}, and the others are not needed for this group
  (item~II.14).
\item The weight $u'_B(k,y)$ is $-H(ky)$, not $|H(ky)|$ (item~II.15,
  Lemma~\ref{lem:uprime}).
\item The invariant spectral terms attached to $\chi\in\mathcal X_A$ do
  not vanish. The logarithmic derivative $m'/m$ of the global normalizing
  factor survives (item~II.16), and $m$ is identified:
  $m'/m=2\frac{L'}{L}(2\lambda,\chi^2)-2\frac{L'}{L}(2\lambda+1,\chi^2)$
  up to finitely many local terms (Corollary~\ref{cor:global-m}).
\item The residual terms $\operatorname{tr}(R(f)|_{V_b})$ and the ramified
  terms of Theorem~\ref{thm:Jchi}(c), which the 1995 notes do not have,
  are already invariant. The ramified term is made explicit: $M(w,0)=-1$
  on $\mathcal H_B(\chi_{\psi_b})$, so it equals
  $-\frac14\operatorname{tr}I(\chi_{\psi_b,0},f)$
  (Proposition~\ref{prop:Mw0}).
\item The function $\varphi_A(f)$ that carries the correction is not the
  Fourier transform of a compactly supported function on $\AS$
  (Example~\ref{ex:Phi-real}); the 1995 definition~(6.7) of the space
  that contains it is circular (item~II.17), and its claim that this
  space equals the Paley--Wiener space of $\AS$ is false.
\end{itemize}

\item[Lemma~\ref{lem:uprime}.] Example~6.1 of the 1995 notes computes $u'_B(k,y)=|H_B(ky)|$ (item~II.15); the correct value is $-H(ky)$.

\item[Theorem~\ref{thm:normalization}.] The theorem replaces Conjecture~6.10 of the 1995 notes; property~(1) there, $R_{P'|P}(\chi)I_P(\chi,x)=I_P(\chi,x)R_{P'|P}(\chi)$, should have $I_{P'}$ on the right, as in~\ref{N1}. At the real places the 1995 notes refer to Sally~\cite{Sa}, who treats the universal cover of $SL(2,\R)$.

\item[Corollary~\ref{cor:global-m}.] Part~(a) is the relation $M^{-1}M'=m^{-1}m'+R^{-1}R'$ of Sections~6.1.1
and~6.2 of the 1995 notes, where the factor $m$ of $R_{\bar B|B}(\sigma)=
m(\sigma)^{-1}M_{\bar B|B}(\sigma)$ is not identified; (c) identifies
it.

\item[The map $\varphi_A$, \eqref{eq:phiA}.] The 1995 notes have $R_A(\pi_\lambda,P_0)=+R^{-1}\frac{d}{d\lambda}R$, without the minus sign.

\item[Theorem~\ref{thm:IA}.] The step that needs the trace Paley--Wiener theorem for a Levi subgroup $L$ is where Conjectures~6.3--6.8 of the 1995 notes enter; for $L=A$ they are not needed. The 1995 condition~(6.7) defining the space that contains $\varphi_A(f)$ is circular: it has $\phi(\sigma,X)$ on both sides (item~II.17). The 1995 notes assert that $\mathcal I_{\mathrm{ac}}(\overline{A(\A_S)})=\mathcal I(\overline{A(\A_S)})$, the space of Fourier transforms of compactly supported functions; this is false, because $\mathfrak a_A\ne0$ (Example~\ref{ex:Phi-real}). The vanishing of $I_A(\pi,X,f)$ is Proposition~6.20 and Example~6.21 there.

\item[Theorem~\ref{thm:itf}.] In Section~6.4 of the 1995 notes, (a) has $\frac12a^A(S,z)I_A(z,f)$ and the sum
over $(G,S)$-classes of $a^G(S,\gamma)I_G(\gamma,f)$, as here
($I_G=J_G$), and (b), (c) agree with ours. The spectral item~(b) there,
$\frac12\int a^A(\pi)I_A(\pi,f)\,d\pi$ with the remark that $I_A(\pi,f)=0$
for tempered $\pi$, is identically zero (item~II.16). The book version
(\texttt{book7g.tex}, item~(e)) has $\int a^G(\pi)\operatorname{tr}\pi(f)\,
d\pi$ with $a^G(I_A(\sigma))=-m'(\sigma)/m(\sigma)$, which is~(e) above;
neither has the residual and ramified terms of~(f). Theorem~6.22 there
(Arthur's~\cite[Cor.~4.5]{A10}; $v_\pi$ should be $\nu_\pi$) is the
statement that $I_\chi(f)=\int a^G(\pi)I_G(\pi,f)\,d\pi$, the term $M=A$
being $0$.

\end{description}


\subsection{Section~\ref{sec:stable}: the stable(r) form}

Section~\ref{sec:stable} replaces Section~7 of the 1995 notes
(\texttt{cctrace7.tex}, which is \texttt{book7h.tex}) and uses the parts
of \texttt{book3b.tex} (``Local comparisons on the groups'') that concern
conjugation by $GL(2)$. The 1995 section was a program modeled
on~\cite{LL}, with most proofs referred to preprints that were not
written; the items below say what became of each part.

\begin{description}
\item[Stable conjugacy (Section~\ref{sec:stable-conj}).] The 1995 notes call two elements of the cover stably conjugate if they are conjugate by an element of $G_1(F)=GL(2,F)$, and a distribution stably invariant if $T(f^y)=T(f)$ for all $y\in G_1(\A_S)$. The cover of $G_1(F)$ is said there to be ``defined using the cocycle $\beta$'', which is defined only on $SL(2)$ (item~II.18); Section~\ref{sec:gl2} fixes one. With that definition the two lifts of $\gamma$ are stably conjugate whenever $\eta_\gamma$ is not trivial on the norms (Lemma~\ref{lem:eta}), every stably invariant distribution vanishes on genuine functions supported near such elements, and the functions $\delta\mapsto I_{\mathcal O}(f^\delta)$ of~(7.5)--(7.8) are defined on $D(T/F)=T_1(F)\backslash G_1(F)/G(F)$ only up to the character $\eta_\gamma$. The action $\star$ of~\eqref{eq:star} removes the ambiguity.

\item[Endoscopic groups.] Section~7.2 of the 1995 notes takes the Cartan subgroups of $SL(2)$ as the endoscopic groups of the cover, as for $SL(2)$ in~\cite{LL}, and writes the ``endoscopic terms''~(7.2) as $\frac14\sum_H\sum_\theta\langle\theta,f^H\rangle$ over characters $\theta$ of the elliptic tori $H$. For the double cover the elliptic endoscopic groups are two copies of $SO(3)$~\cite{Li11}, and the $\kappa$-orbital integrals are not functions to which Poisson summation on the torus is applied: they are stable orbital integrals of $f_-$ (Lemma~\ref{lem:kappa-orbital}(c)). Accordingly Proposition~7.2 (Poisson summation on $T$), Corollary~7.4, Lemma~7.20, Lemma~7.21 and the sums~(7.20)--(7.21) over characters $\theta$ of tori have no counterpart here. In \texttt{book3b.tex} the $L$-group of the cover is said to be $PGL(2,\C)$ (item~C.4); it is $SL(2,\C)$ (Section~\ref{sec:stable-conj}).

\item[Theorem~\ref{thm:elliptic-stable}.] This replaces (7.3)--(7.8), Lemma~7.5 and the elliptic part of~(7.20). Three corrections (item~II.19): $T(\Q)$ is not discrete in $T(\A_S)$, so the statement is adelic; the index $[D(T/\Q):D(T/\A_S)]^{-1}$ is the factor $\frac12$, by the Hasse norm theorem and class field theory; and Lemma~7.5(a), as stated, is about $I^{\mathrm{st}}_{\mathcal O}$, not $SI_{\mathcal O}$. The local sets $D(T/F)$ are $F^\times/N E^\times$.

\item[Proposition~\ref{prop:IA-gl2} and Lemma~\ref{lem:variance-gl2}.] Section~7.4 of the 1995 notes forms $D(A/F)$ with $A\cdot Z$ in place of the centralizer of $A$ in $G_1$, which gives $F^\times/F^{\times2}$, and sums over it with an index $[D(A/\Q):D(A/\A_S)]$ that is not defined, since $\Q^\times/\Q^{\times2}$ does not embed in a finite group (item~II.20). For a hyperbolic element the stable class is the conjugacy class; the group $F^\times/F^{\times2}$ acts on its two lifts through $(a,t)_v$ (Lemma~\ref{lem:Ad}(d)), and Proposition~\ref{prop:IA-gl2}(b) is what remains of~(7.9)--(7.10). Lemma~7.6 there (stability of $SI_{\mathcal O}$ for these classes) is not proved and is not claimed here (see the remark after Remark~\ref{rem:N5}). Lemma~7.7 is the identity $\operatorname{tr}I(\pi,f)=\int\pi(\tilde a)J_G(\tilde a,f)\,d\tilde a$ used in Proposition~\ref{prop:IA-gl2}(c). Lemmas~7.8, 7.9 and~7.23, on the operator $N(g)$ of multiplication by $\log\delta(kg)$ and the variance of $J_A$ and of $\operatorname{tr}(R^{-1}R'I_A(\sigma,f))$ under $g\in G_1(F)$, are Lemma~\ref{lem:variance-gl2}; the 1995 statements omit the change $\gamma\mapsto\mathrm{Ad}(\delta)\gamma$, $\sigma\mapsto\sigma\kappa_a$ and the condition~\ref{N5}. In the proof sketched in \texttt{book3b.tex} the matrix identity for $wn(x)$ is wrong (item~C.1); the correct factorization is $wn(x)=n(-x^{-1})\,a(x^{-1})\,\bar n(x^{-1})$, and the weight is $v_A=-H(x)-H(wx)$ (Lemma~\ref{lem:height-ineq}). The identification of $K$ with the subset $(K,1)$ of the cover in Section~7.4 (item~I.1) is not used.

\item[The weight~(7.13).] The formula $v_A=\sum_v\log\sqrt{|1+n_v^2|_v}$ holds only at real places (item~II.22); see Lemma~\ref{lem:weight-local}.

\item[Proposition~\ref{prop:unipotent}.] This replaces Section~7.5 of the 1995 notes. Equation~(7.17) there sums over the rational unipotent classes $a\in\Q^\times/\Q^{\times2}$ and diverges, and its middle equality is wrong (item~II.21); the correct expansion is over the finitely many quadratic Hecke characters unramified outside $S$, with the coefficients $L^S(1,\kappa)$. Definitions~7.12 and~7.13 use the transfer factor $\lambda(E/F,\psi)\kappa\bigl((\gamma_1-\gamma_2)/(\gamma_1^0-\gamma_2^0)\bigr)$ of~\cite{LL} for $SL(2)$; for the cover the corresponding function is $\Delta_\psi$ of Lemma~\ref{lem:Delta}. Lemmas~7.14--7.16 (smoothness of the normalized $\kappa$-orbital integrals on the torus and their limits at $\pm1$) are not proved, here or there; the statement of Lemma~7.15 lacks an equality sign and has a constant depending on $\gamma$ inside a limit.

\item[Proposition~\ref{prop:theta-packet} and Lemma~\ref{lem:spectrum-gl2}.] Definition~7.18 of the 1995 notes introduces the set $L(\pi)$ of the $\pi\circ\mathrm{Ad}(g)$, $g\in G_1(F)$, and Remark~7.19 asserts that it has $\prod_{v\mid p}\#(k_v^\times/k_v^{\times2})$ elements for every genuine $\pi$ (item~II.23; item~C.2). This is not so in general: $I(\pi)\circ\mathrm{Ad}(\delta_a)=I(\pi\kappa_a)$ (Proposition~\ref{prop:IA-gl2}(c)), which is equivalent to $I(\pi)$ when $\pi\kappa_a=\pi^w$, that is, when $\chi_v^2=(a,\cdot)_v$; and $r^\pm_{\psi}\circ\mathrm{Ad}(\delta_a)=r^\pm_{\psi_a}$ (Lemma~\ref{lem:weil-kernel}). The sets $L(\pi)$ are not the packets of the cover; those are Waldspurger's, of which Proposition~\ref{prop:theta-packet} treats one family. Lemma~7.20 states no assertion.

\item[Theorem~\ref{thm:stabler}.] This replaces~(7.1), (7.25) and the list (a)--(i) of Section~7.8. The distributions $S_1$, $S_2$, $S_3$ of~(7.22)--(7.24) are, in the present notation, the correction $\Phi_A$, the weighted orbital integrals $J_A$ (the two together are the terms $I_A$) and the unipotent terms $J^\kappa_{\mathrm{unip}}$. Conjecture~7.24, that $S_1+S_2+S_3$ is invariant under $G_1(\A_S)$, does not hold for the cover as stated: by Proposition~\ref{prop:IA-gl2}(b) the term of $a(t)$ is multiplied by $(\det h,t)_S$. What replaces it is the equivariance in Theorem~\ref{thm:stabler}(c), and Li's stabilization~\cite{Li21}. The term~(g) of Section~7.8, with $m'/m$, is Theorem~\ref{thm:itf}(e)--(f).

\item[\texttt{book3b.tex}.] Its first lemmas assert that the orbital integral of a genuine function at a lift of $a\in A(F)$ vanishes unless $a$ is an $N$-th power. For $n=2$ the centralizer of the lift is the full preimage of $A$ (Lemma~\ref{lem:classes}) and the orbital integral does not vanish; what vanishes for non-squares is its average over the $G_1(F)$-orbit (Lemma~\ref{lem:kappa-orbital}(b)). The correspondence $\chi\leftrightarrow\bar\chi$ between characters of $A(F)$ and genuine characters of its cover, defined there by cases, is multiplication by Weil's character $\chi_\psi$ (Lemma~\ref{lem:transfer}), whose inverse is the restriction of $\Delta_\psi$ to the cover of $A$ (Lemma~\ref{lem:Delta}(c)). The real Hilbert symbol is misdefined there as the sign of the first argument (item~C.3); see~\ref{H3}.

\end{description}
